\chapter{Descrizione dello stage}
\label{chap:descrizione-stage}

\textit{Il presente capitolo descrive il percorso di stage svolto presso l'azienda, illustrando gli obiettivi formativi,
le competenze richieste e le attività pianificate durante il periodo di tirocinio. Vengono inoltre presentate
l'organizzazione del lavoro adottata, l'analisi preventiva dei principali rischi progettuali e il ruolo svolto
dal tutor aziendale nel supporto alle attività di sviluppo.}

\section{Competenze e obiettivi}

L'obiettivo principale dello stage è stato l'evoluzione e l'ottimizzazione del firmware embedded del sistema
\textit{EggLog Motion Core}, una piattaforma dedicata all'acquisizione e alla gestione di dati telemetrici
in ambito motorsport, sviluppata su piattaforma ESP32.

Il lavoro riguarda principalmente l'analisi dell'architettura firmware esistente, il miglioramento dei
componenti software già presenti e l'introduzione di nuove funzionalità volte a rendere il sistema più
affidabile, efficiente e facilmente estendibile.

Le attività svolte hanno coinvolto diversi aspetti del sistema, tra cui la gestione dello \textit{storage}
su memoria SD, l'ottimizzazione dei servizi di comunicazione BLE (\gls{bleg}) e Wi-Fi, la gestione del
trasferimento dei dati acquisiti e l'analisi del traffico CAN proveniente dalla centralina elettronica
(\gls{ecug} detta ECU) della moto mediante tecniche di \gls{cansniffing}.

\subsection{Competenze da acquisire}

Durante il periodo di tirocinio erano previste l'acquisizione e il consolidamento delle seguenti competenze:

\begin{itemize}

    \item sviluppo e manutenzione di firmware embedded in linguaggio \gls{cpp} su piattaforma ESP32;

    \item analisi, \textit{refactoring} e ottimizzazione di codice embedded esistente;

    \item gestione di sistemi di memorizzazione embedded e ottimizzazione delle operazioni di lettura e scrittura
    su memoria SD;

    \item sviluppo e ottimizzazione di servizi \gls{ble} e Wi-Fi;

    \item analisi dei protocolli di comunicazione utilizzati in ambito automotive, con particolare riferimento
    al \gls{can};

    \item acquisizione e analisi di messaggi dal \gls{can} provenienti dalla \gls{ecu} di un veicolo motorizzato;

    \item utilizzo di strumenti di \textit{debugging} e \textit{testing} per sistemi embedded;

    \item produzione di documentazione tecnica relativa al firmware e ai protocolli di comunicazione analizzati.

\end{itemize}


\subsection{Obiettivi}

Gli obiettivi definiti dall'azienda sono stati classificati secondo tre livelli di priorità:

\begin{itemize}
    \item \textbf{O}: obbligatorio;
    \item \textbf{D}: desiderabile;
    \item \textbf{F}: facoltativo.
\end{itemize}


\subsubsection{Obiettivi obbligatori}

Gli obiettivi obbligatori comprendevano principalmente attività volte al miglioramento dell'affidabilità
e delle prestazioni del firmware esistente.

\begin{itemize}
    \item \textbf{OO-1}: \textit{Refactoring} e ottimizzazione del sistema di lettura e scrittura su scheda SD;

    \item \textbf{OO-2}: Completamento del sistema di rotazione dei file con introduzione di meccanismi
    di recupero in caso di interruzioni impreviste;

    \item \textbf{OO-3}: Ottimizzazione delle procedure di download dei dati memorizzati sulla scheda SD;

    \item \textbf{OO-4}: Miglioramento dei servizi \gls{ble} e introduzione di nuovi comandi firmware;

    \item \textbf{OO-5}: Produzione della documentazione tecnica relativa al firmware sviluppato.
\end{itemize}


\subsubsection{Obiettivi desiderabili}

Gli obiettivi desiderabili erano principalmente orientati all'introduzione di funzionalità aggiuntive
e al miglioramento delle modalità di trasferimento dei dati.

\begin{itemize}
    \item \textbf{OD-1}: Implementazione del trasferimento dei dati tramite connessione Wi-Fi;

    \item \textbf{OD-2}: Introduzione di meccanismi di supporto alla ripresa dei trasferimenti interrotti;

    \item \textbf{OD-3}: Valutazione e introduzione di tecniche di compressione dei dati;

    \item \textbf{OD-4}: Integrazione dell'analisi e gestione dei dati proveninti dalla \gls{ecu};

    \item \textbf{OD-5}: Raccolta di metriche relative alle prestazioni dei trasferimenti dati con 
    relativi \gls{benchmark} per il confronto.
\end{itemize}


\subsubsection{Obiettivi facoltativi}

\begin{itemize}
    \item \textbf{OF-1}: Implementazione di un sistema avanzato di indicizzazione delle sessioni
    memorizzate sulla scheda SD.
\end{itemize}


\section{Pianificazione}

Le attività di stage sono state pianificate con il tutor aziendale considerando un monte complessivo
di 320 ore, distribuite su otto settimane lavorative da quaranta ore ciascuna. Tali settimane sono state 
suddivise come segue:

\begin{enumerate}

    \item \textbf{Prima settimana -- \textit{Onboarding} e analisi del progetto}
    \begin{itemize}
        \item Introduzione al sistema \textbf{EggLog Motion Core} e all'architettura generale della piattaforma;
        \item Configurazione dell'ambiente di sviluppo e degli strumenti di compilazione e debugging;
        \item Analisi della struttura del firmware esistente e delle principali componenti software;
        \item Studio delle modalità di acquisizione, memorizzazione e trasferimento dei dati.
    \end{itemize}

    \item \textbf{Seconda settimana -- Studio delle tecnologie e analisi del firmware}
    \begin{itemize}
        \item Approfondimento della gestione della memoria microSD e del relativo \textit{file system};
        \item Analisi dello stack \gls{ble} e dei servizi implementati;
        \item Studio dei protocolli di comunicazione utilizzati dal sistema;
        \item Individuazione delle principali criticità presenti nel firmware.
    \end{itemize}

    \item \textbf{Terza settimana -- Analisi e progettazione delle modifiche}
    \begin{itemize}
        \item Analisi dell'architettura software esistente e definizione degli interventi di \textit{refactoring};
        \item Progettazione delle modifiche relative alla gestione dello storage;
        \item Studio dei meccanismi di recupero e gestione dei dati in caso di errore;
        \item Analisi preliminare del \gls{can} della moto e delle informazioni trasmesse dalla \gls{ecu}.
    \end{itemize}

    \item \textbf{Quarta settimana -- Sviluppo e ottimizzazione dello storage}
    \begin{itemize}
        \item \textit{Refactoring} dei moduli responsabili della gestione della memoria SD;
        \item Ottimizzazione delle operazioni di lettura e scrittura dei dati;
        \item Miglioramento del sistema di rotazione dei file;
        \item Implementazione e verifica dei meccanismi di \textit{recovery}.
    \end{itemize}

    \item \textbf{Quinta settimana -- Gestione dei dati e comunicazioni}
    \begin{itemize}
        \item Implementazione delle funzioni di calibrazione dei sensori;
        \item Ottimizzazione delle procedure di download dei dati memorizzati;
        \item Studio delle possibilità di compressione dei dati;
        \item Miglioramento dei servizi \gls{ble}.
    \end{itemize}

    \item \textbf{Sesta settimana -- Analisi del \gls{can} e integrazione firmware}
    \begin{itemize}
        \item Acquisizione dei messaggi presenti sul \gls{can} tramite attività di \textit{sniffing};
        \item Analisi degli identificativi e dei dati trasmessi nel \gls{can} dalla \gls{ecu};
        \item Studio della possibile integrazione delle informazioni trasmesse nel \gls{can} all'interno 
        del sistema telemetrico;
        \item Evoluzione dei comandi e dei servizi firmware.
    \end{itemize}

    \item \textbf{Settima settimana -- Testing e validazione}
    \begin{itemize}
        \item Esecuzione di test funzionali sul firmware modificato;
        \item Raccolta di metriche sulle prestazioni del sistema;
        \item Esecuzione di benchmark sulle operazioni di memorizzazione e trasferimento dati;
        \item Analisi e risoluzione delle problematiche emerse durante i test.
    \end{itemize}

    \item \textbf{Ottava settimana -- Documentazione finale}
    \begin{itemize}
        \item Aggiornamento della documentazione tecnica del firmware;
        \item Descrizione delle modifiche architetturali introdotte;
        \item Documentazione relativa ai protocolli di comunicazione analizzati;
        \item Stesura della relazione finale di stage.
    \end{itemize}

\end{enumerate}


\section{Organizzazione del lavoro}

Le attività di sviluppo sono state organizzate seguendo una metodologia \gls{agile}, ispirata al framework
\gls{scrum}. Il lavoro è stato suddiviso in sprint della durata indicativa di due settimane, durante i quali
venivano pianificate le attività da svolgere, definiti gli obiettivi da raggiungere e monitorato lo stato
di avanzamento delle attività insieme al tutor aziendale.
\\
Durante il periodo di stage sono stati inoltre svolti degli \gls{standupmeeting} con cadenza giornaliera,
finalizzati al mantenimento dell'allineamento tra i membri del team. Questi incontri hanno permesso di
condividere i progressi effettuati, discutere eventuali problematiche emerse durante le fasi di sviluppo
e progettazione e individuare possibili soluzioni attraverso il confronto tecnico.
\\

Il gruppo di lavoro era composto da due tirocinanti, ciascuno responsabile dello sviluppo di una specifica
componente del sistema. Il sottoscritto ha operato principalmente sulla componente firmware embedded,
occupandosi dell'evoluzione del software a bordo del dispositivo, mentre la collega ha lavorato sulle
componenti applicative del sistema.
\\
Nonostante la suddivisione delle responsabilità, le attività sono state svolte in maniera collaborativa,
favorendo il confronto continuo tra i membri del team e garantendo una corretta integrazione tra firmware,
applicazione mobile e infrastruttura software.

\section{Analisi preventiva dei rischi}
Prima dell'avvio delle attività è stata effettuata un'analisi preventiva dei principali rischi
che avrebbero potuto influenzare il corretto svolgimento del progetto.

\begin{itemize}

    \item \textbf{Comprensione del firmware esistente}\\
    la necessità di intervenire su un sistema software già sviluppato ha richiesto una fase iniziale
    di analisi dell'architettura firmware e dei meccanismi implementati, al fine di ridurre il rischio
    di introdurre regressioni durante le modifiche.

    \item \textbf{Gestione dello storage e integrità dei dati}\\
    le operazioni di lettura, scrittura e trasferimento dei dati dalla memoria microSD rappresentavano una
    possibile criticità in termini di affidabilità. Il rischio è stato mitigato attraverso attività di
    \textit{testing}, analisi delle prestazioni e utilizzo di \gls{benchmark}.

    \item \textbf{Analisi del traffico \gls{can}}\\
    l'acquisizione e l'interpretazione dei dati provenienti dalla \gls{ecu} della moto rappresentavano
    una possibile criticità, in quanto i messaggi presenti sul \gls{can} non seguono necessariamente
    una codifica standardizzata.\\
    Inoltre, prima dell'attività di \textit{sniffing}, è stata quindi svolta un'analisi
    di fattibilità sia dal punto di vista hardware, valutando l'interfacciamento fisico con il veicolo,
    sia dal punto di vista software, considerando l'acquisizione, l'elaborazione e l'interpretazione
    dei dati raccolti.

    \item \textbf{Coordinamento tra le componenti del sistema}\\
    la presenza di più componenti sviluppate parallelamente poteva introdurre problematiche di
    integrazione tra firmware, applicazione mobile e infrastruttura software. Tale rischio è stato
    mitigato attraverso confronti periodici con il tutor aziendale, revisioni del codice e attività
    di integrazione condivise.

\end{itemize}